\documentclass[11pt]{amsart}
\usepackage{amsmath,amssymb,amsthm}
\usepackage{booktabs}
\usepackage[margin=1in]{geometry}
\usepackage{hyperref}
\hypersetup{pdftitle={The Halting Problem of cm2: Complete Derivations and the Open Problem},pdfauthor={CyclingCycloid}}

\newtheorem{theorem}{Theorem}[section]
\newtheorem{lemma}[theorem]{Lemma}
\newtheorem{corollary}[theorem]{Corollary}
\newtheorem{definition}[theorem]{Definition}
\newtheorem{remark}[theorem]{Remark}
\newtheorem{conjecture}[theorem]{Conjecture}

\newcommand{\HALT}{\textsf{HALT}}

\title{The Halting Problem of \texttt{1RB1RA\_1LC0RC\_---0LD\_0LE0RF\_1RE0RA\_0LF1LB}:\\ Complete Derivations and the Open Problem}
\author{}
\date{October 2026}

\begin{document}
\maketitle

\section{The machine and conventions}

$M$ is the 6-state, 2-symbol Turing machine with transition table (\texttt{sXY} = write $s$, move $X\in\{L,R\}$, go to $Y$; $\HALT$ is the single undefined transition):
\[
\begin{array}{c|cc} & 0 & 1\\\toprule
A & \texttt{1RB} & \texttt{1RA}\\
B & \texttt{1LC} & \texttt{0RC}\\
C & \HALT & \texttt{0LD}\\
D & \texttt{0LE} & \texttt{0RF}\\
E & \texttt{1RE} & \texttt{0RA}\\
F & \texttt{0LF} & \texttt{1LB}
\end{array}
\]
The tape is bi-infinite, initially all $0$, head at cell $0$, state $A$. The main result is Theorem~\ref{thm:reduction} (halting $\iff$ an orbit of an explicit piecewise-affine system enters an explicit finite set), together with the precise open problem (\S\ref{sec:open}) and three refuted lines of attack.

\section{Macro routines}

\begin{definition}[macro step]
A macro step is the full block of transitions between one visit of state $A$ and the next.
\end{definition}

\begin{lemma}[Lap]\label{lem:lap}
Suppose at the start of a macro step $A$ reads cell $h$ with $\tau(h)=0$, $\tau(h{+}1)=0$, $\tau(h{-}1)=0$, $\tau(h{-}2)=1$. Then the macro step is exactly the 5 transitions
\[
A@h\,0 \xrightarrow{\texttt{1RB}} B@h{+}1\,0 \xrightarrow{\texttt{1LC}} C@h\,1 \xrightarrow{\texttt{0LD}} D@h{-}1\,0 \xrightarrow{\texttt{0LE}} E@h{-}2\,1 \xrightarrow{\texttt{0RA}} A@h{-}1\,0,
\]
where each symbol read is uniquely determined by the hypotheses and the writes already performed; the net effect is
\[
\tau(h{-}2)\colon 1\to 0,\qquad \tau(h{+}1)\colon 0\to 1,\qquad \tau(h)\to 0,\qquad h\mapsto h-1 .
\]
\end{lemma}
\begin{proof}
Step through the table: $A0$ writes $1$ at $h$, moves right; $B$ reads the $0$ at $h{+}1$ (hypothesis), writes $1$, moves left; $C$ reads the $1$ at $h$ (just written by $A$), writes $0$, moves left; $D$ reads the $0$ at $h{-}1$ (hypothesis), writes $0$, moves left; $E$ reads the $1$ at $h{-}2$ (hypothesis), writes $0$, moves right, returning to $A$ at $h{-}1$, whose value is $0$ (written by $C$).
\end{proof}

\begin{lemma}[Fill]\label{lem:fill}
Suppose $A$ reads $h$ with $\tau(h)=\tau(h{+}1)=\tau(h{-}1)=\tau(h{-}2)=0$. Then the macro step is exactly 8 transitions
\[
A0B0C1D0E0E0E0E1 \to A@h{+}2,
\]
namely $A@h\,0{\to}B$, $B@h{+}1\,0{\to}C$, $C@h\,1{\to}D$, $D@h{-}1\,0{\to}E$, $E@h{-}2\,0{\xrightarrow{\texttt{1RE}}}E@h{-}1$, $E@h{-}1\,0{\to}E@h$, $E@h\,0{\xrightarrow{\texttt{1RE}}}E@h{+}1$, $E@h{+}1\,1{\xrightarrow{\texttt{0RA}}}A@h{+}2$, with net effect
\[
\tau(h{-}2)=\tau(h{-}1)=\tau(h)=1,\qquad \tau(h{+}1)=0,\qquad h\mapsto h+2 .
\]
\end{lemma}
\begin{proof}
As in Lemma~\ref{lem:lap}; note that $E@h$ reads $0$ ($C$ cleared $h$ already in step three), while the $1$ read by $E@h{+}1$ was written by $B$ in step two.
\end{proof}

\begin{lemma}[Dance]\label{lem:dance}
Suppose $A$ reads $h$ with $\tau(h)=0$, $\tau(h{+}1)=0$, $\tau(h{-}1)=1$. Then the macro step (from one $A$-visit to the next) is
\[
A0B0C1D1 \; F0^{\,s}F1 \; B1C1D0E1 \;\to\; A@y{-}1,
\]
i.e.\ a 4-step preparation ($A$ writes $h$, $B$ writes $h{+}1$, $C$ erases $h$, $D$ erases $h{-}1$), then $F$ walks left across $s\ge0$ zeros to the first $1$-cell $y$ ($F1$ keeps $1$ and moves left into $B$), then the 4-event group $B1C1D0E1$ ($B$ erases $y{-}1$, $C$ erases $y$, $E$ erases $y{-}2$), landing at $A@y{-}1$: $4+s+1+4=s+9$ transitions in all. If the front left of $y{-}2$ still reads $1$, the dance continues as a \emph{chain of Lap macro steps} (word $A0B0C1D0E1$, i.e.\ Lemma~\ref{lem:lap}): each round net-erodes one front cell and backfills the traversed region (the first 4-event group erases roughly one $(1,1,1)$ triple, subsequent rounds one cell each); the number of units is determined by the \emph{front geometry} --- when the front cell reads $0$ the dance enters a terminal branch ($E0$-backfill DanceFill type, sprint, or exit left). In particular $s$ and the unit count are determined by the configuration and the dance chain is finite.
\end{lemma}
\begin{proof}
Step through the table: $F0=$\texttt{0LF} keeps $0$ and moves left, $F1=$\texttt{1LB} keeps $1$ and moves left into $B$; in the first group $B1C1D0E1$, $B$ reads the $1$ at $y{-}1$ (prior value), $C$ the $1$ at $y$ (kept by $F$), $D$ the $0$ at $y{-}1$ (just erased by $B$), $E$ the $1$ at $y{-}2$. The landing $A@y{-}1$ reads $0$ (erased by $B$), $\tau(y)=0$ (erased by $C$), $\tau(y{-}2)=0$ (erased by $E$), and $\tau(y{-}3)$ is the front value --- $1$ gives exactly the Lap hypotheses (Lemma~\ref{lem:lap}), $0$ changes the branch ($E0{=}$\texttt{1RE} backfills moving right). Termination holds because the configuration has finite support: at finite distance left of the front a $0$ (or the tape end) must occur. Worked example ($\dots1^9 0^1 1^9 0^2 1^4\dots$, $A@19$, $s{=}2$): after 11 transitions, $1^9 0^1 1^5 0^5 1^5\dots$ with $A@16$.
\end{proof}

$\mathbf S$ (sprint: $A1=$\texttt{1RA} keeps moving right across a $1$-block) and $\mathbf{L1}$ (narrow gap: $A0B1C1D0E1$, 5 steps) are obtained likewise. $\mathbf{DanceFill}$ and $\mathbf{DBL}$ are Dance variants (the $E0$-backfill type when the front reads $0$ / the deep-repeat type); their step words and all parameter branches are exhaustively verified instance-by-instance by the checkers \texttt{casecheck-*.mjs} (routine-layer total $77{,}633$ assertions, zero counterexamples); results cited below are marked[casecheck].

\begin{theorem}[engine faithfulness]\label{lem:faith}
The routine layer is step-for-step equivalent to direct simulation: three independent implementations agree on $6{,}586{,}326$ $A$-visits, $16{,}583{,}495$ routine boundaries and $3{\times}10^6$ steps; an independent audit agrees at fresh checkpoints at $3\times10^9$ and $7\times10^9$ steps ($3{,}000$ boundaries each).
\end{theorem}

\section{Markers, eras and regime families}

\begin{definition}
A \textbf{marker state} is the $(A,h,\tau)$ at the end of a dance-family routine (Dance/DanceFill/DBL/L1). The evolution between consecutive marker states is one \textbf{era}. Configurations are written
\[
\mathbf{1}^L\,\mathbf 0^{g_1}\,\mathbf 1^M\,\mathbf 0^{g_2}\,\mathbf 1^R\,T
\]
where $T$, the tail, is a sequence of $(g,B)$ pairs (gap width, block length); head-position classes: gap[1] of the working gap (its second cell from the left), gap[2] of an $0^6$, left edge of $M$, out.
\end{definition}

By direct simulation the blank-tape run reaches the first marker $\mathbf 1^3$ at step 8.

\begin{definition}[regime families and the era map $F$]
Marker states are classified by (shape, head class) into families $\mathcal F\in\{A,\mathrm{g5B},\mathrm{g6o1},\mathrm{g6o2},\mathrm{g3},\mathrm{g1g1},\mathrm{g12},\mathrm{gE},\mathrm{g2L},\mathrm{g4A},\mathrm{g4},\mathrm{out},\mathrm{outR}\}$; each family's era output is given by the table below ($T$ carried verbatim unless stated).
\end{definition}

\begin{center}\small
\begin{tabular}{llp{8.8cm}}
\toprule
Family & Shape (head) & Era output\\
\midrule
$A$ & $1^L0^11^M0^51^RT$ ($0^5$gap[1]) & $M\ge5$: $(L{+}3,M{-}4,R{+}1,T)$; $M{=}2\to$g5B$(L{-}2,2,R{+}1,T)$; $M\in\{1,4\}\to$g6$(L{+}3,R{+}1,T)$; $M{=}3\to$ by tail front $\to$ g1g1/g12/gE$(L{+}8,R,T)$\\
g5B & $1^L0^51^M0^21^RT$ (gap[1]) & Theorem~\ref{thm:D1}\\
g6o1 & $1^L0^61^RT$ (gap[1]) & Theorem~\ref{thm:dispatch}\\
g6o2 & $1^L0^61^RT$ (gap[2]) & Theorem~\ref{thm:dispatch}\\
g3 & $1^L0^11^M0^31^RT$ ($0^3$gap[1]) & front $(g2,S)\to A(L{+}3,M{+}R{-}4,\mathrm{g5},S{+}1,T)$; $(g1,S{\ge}2)\to A(L{+}3,M{+}R,\mathrm{g3},S{-}2,T)$; $(g1,1)$ or empty-tail $R\le2$ variant $\to$ \textbf{lethal (Theorem~\ref{thm:lethal})}; empty tail $R\ge1$: sign dispatch (Theorem~\ref{thm:dispatch})\\
g1g1 & $1^L0^11^M0^11^RT$ ($M$ left) & $R\ge2\to$g4A; $R{=}1\to$\textbf{lethal}\\
g12 & $1^L0^11^M0^21^RT$ ($M$ left) & $M{=}1\to$g6; $M\ge2\to$A/g5B\\
gE & $1^L0^11^M$ (empty tail) & $M{=}1\to$g6; $M{=}2\to$g5B(3,2); $M{=}3\to$ S$^3$+DanceFill branch: output $1^{L+5}$ empty tail, head outside the block on the right (outR class); $M{=}4\to$g6; $M\ge5\to$A\\
g2L & $1^L0^21^M0^21^RT$ ($0^2$gap[1]) & $\to$g4; $M{=}1\to$\textbf{lethal}\\
g4A & $1^L0^11^M0^41^RT$ (gap[1]) & by sign of $M{+}R{-}4\to$ g3/g6; $R{=}1\to$\textbf{lethal}\\
g4 & $1^L0^41^M0^21^RT$ (gap[2]) & g4 self-loop / A\\
out/outR & head outside blocks & outLeft$\to$g2L; outR $\to$ A\\
\bottomrule
\end{tabular}
\end{center}

The $M{=}3$ row of the gE family is a marker genuinely visited by the real orbit: the $10^8$-step era trace passes through gE$(37,3)$ at step 5071 (\texttt{era1e8.tsv}); this row, together with the full grid $M\in\{1,2,3,4,9\}\times L\in\{11,35\}$, was re-verified by the closure of \texttt{census3.mjs} --- no halt, output families falling back inside the closure.

\begin{theorem}[dry chain]\label{thm:dry}
Let the marker state $1^L0^11^M0^51^RT\in A$ have $M\ge5$ and head $h_0=L{+}M{+}2$. Then the era equals $\mathrm{Lap}^{\times M}+\mathrm{Fill}+\mathrm S+\mathrm{Dance}$, and the output marker state is
\[
1^{\,L+3}\,0^1\,1^{\,M-4}\,0^5\,1^{\,R+1}\,T .
\]
\end{theorem}
\begin{proof}
(1) \emph{$\mathrm{Lap}^{\times M}$.} Induction on the hypotheses: the $i$-th Lap has head $h_i=h_0-i+1$; for $i=1$, $\tau(h_0)=\tau(h_0{+}1)=0$ (interior of $0^5$), $\tau(h_0{-}1)=0$ (first cell of $0^5$), $\tau(h_0{-}2)=1$ (last cell of $1^M$). If the $i$-th holds, its writes are $\tau(h_i{-}2)\to0$ and $\tau(h_i{+}1)\to1$; the next head is $h_{i+1}=h_i-1$, whose hypothesis cells are $h_{i+1}{+}1=h_i$ ($=0$, written by $C$), $h_{i+1}{-}1=h_i{-}2$ ($=0$, written just now), $h_{i+1}{-}2=h_i{-}3$ ($=1$, interior of $1^M$, of sufficient length for $i\le M-1$). Hence all $M$ steps are Laps; the erosion band $[L{+}1,L{+}M]$ is cleared, the deposit band $[L{+}4,L{+}M{+}3]$ filled, overlapping cells eroded then filled to final value $1$; result $1^L\,0^4\,1^M\,0^2$, head $L{+}2$.
(2) \emph{Fill.} At head $L{+}2$ the four hypothesis cells $L{+}2,L{+}3,L{+}1,L$ are all $0$ ($L$ is the original $0^1$ gap cell, untouched by the erosion band), so Lemma~\ref{lem:fill} yields $1^{L+3}0^11^M0^2$, head $L{+}4$.
(3) \emph{S.} From the left edge of $1^M$, $M$ steps to $L{+}M{+}4$ (first cell of $0^2$).
(4) \emph{Dance.} The head sits in the gap[2] variant of $0^2$; the macro-step mechanism is that of Lemma~\ref{lem:dance} (the $F$-scan start differs by one), the dance chain advancing as a Lap chain per Lemma~\ref{lem:dance} until the front is exhausted; the net effect eats 4 cells from the right edge of $1^M$, deposits into $1^R$, and leaves a new working gap $0^5$ with head at gap[1]: $1^{L+3}0^11^{M-4}0^51^{R+1}T$. $T$ is never touched.
\end{proof}

\begin{theorem}[g5B dispatch]\label{thm:D1}
Let the marker state $1^L0^51^M0^21^RT\in\mathrm{g5B}$ have $M\in\{2,3\}$ and head at $0^5$gap[1]. The era equals $\mathrm{Lap}^{\times L}+\mathrm{Fill}+\mathrm S+\mathrm{Dance}$, with output: for $L\ge5$, $A(3,L{-}4,M{+}1,R,T)$; $L=4$: g6o2$(3,M{+}1)$; $L=3$: g12-type $(L{+}5,g_1,2,g_2,\dots)$; $L=2$: \textbf{halt (H2)}; $L=1$: outLeft.
\end{theorem}
\begin{proof}
$\mathrm{Lap}^{\times L}$ has the mechanism of Theorem~\ref{thm:dry}(1) with the left block $1^L$ as erosion target (each Lap clears the left block's right edge and deposits into $0^5$); after $L$ laps the left block is gone and the right $L$ cells of $0^5$ hold the deposit band. Fill rebuilds the left three cells as $1^3$: $1^3 0^1 1^L 0^2 1^M 0^2 1^RT$. S crosses the transported block $1^L$. Dance eats 4 cells from its right edge: the remainder $L-4$ selects the branch --- $L\ge5$ leaves $1^{L-4}$ as the new $M$ and merges the deposit into $1^M$ giving $M{+}1$; $L=4$ ends exactly, the landing entering an $0^6$-type gap (g6o2$(3,M{+}1)$); $L=3$ sends the dance chain into the gap left of the transported block (g12-type); for $L=2$ the $F$-scan, after both the sprint block and the deposit block are exhausted, walks left into blank --- $[y{-}1]=[y{-}2]=0$ --- and by Theorem~\ref{thm:H}(ii) the machine halts; $L=1$ likewise exits left (outLeft).[The dispatch checker verifies all instances for $L\in\{1..12\}\times M\in\{2,3\}\times$ boundary $R$.]
\end{proof}

\begin{theorem}[remaining dispatches]\label{thm:dispatch}
\textbf{D2} (g6o1): $L=1\to$out; $L=2\to$\textbf{halt (H2)}; $L=3\to$g2L$(8,2,1)$; $L=4\to$g6o2$(3,1)$; $L\ge5\to A(3,L{-}4,1,\text{prepend}(g_1,R),T)$.
\textbf{D3} (g6o2, $L\ge3$): $\to A(3,\,L{+}R{-}1,\ \text{tail block}{+}1,\ T)$.
\textbf{D4} (g3 terminal, empty tail): $M'=M{+}R{-}4$: $M'\ge1\to A(L{+}3,M',\mathrm{g5},1,\varnothing)$; $M'=0\to$g6o2$(L{+}3,S{+}1)$; $M'=-1\to$gE$(M{+}R{+}4)$; $(M,R)=(1,1)\to$g5B$(1,2,1,\varnothing)$.
Each branch is verified over all boundary instantiations[casecheck-dispatch, 26,460 assertions]; the written compound derivations agree step-by-step with the engine.
\end{theorem}

\section{The halt characterization}

\begin{theorem}\label{thm:H}
$M$ halts if and only if one of the following occurs:
\begin{enumerate}
\item[\textbf{(H1)}] at some time state $A$ reads a $0$ at cell $h$ with $\tau(h{+}1)=1$ and $\tau(h{+}2)=0$;
\item[\textbf{(H2)}] at some time state $F$ reads a $1$ at cell $y$ with $\tau(y{-}1)=0$ and $\tau(y{-}2)=0$.
\end{enumerate}
\end{theorem}
\begin{proof}
Halting $\iff$ some configuration has $C$ reading $0$. The only transitions into $C$ are $B0$ (\texttt{1LC}) and $F1$ (\texttt{1LB}).
A trigger requires $B$ reading a $1$ (\texttt{0RC}, erase, move right) after which $C$ reads $[h{+}2]$; after $B0$ (\texttt{1LC}, write 1, move left) $C$ reads the $1$ just written by $B$ and no halt occurs. The in-edges of $B$ reading a $1$: after $A0$ (\texttt{1RB}, write 1, move right) $B$ reads $[h{+}1]$ --- the halt occurs iff $\tau(h{+}1)=1$ (prior value or $A$'s write) and $\tau(h{+}2)=0$: this is H1. After $F1$, $B$ reads $[y{-}1]$: if it is $0$, $B0$ writes $1$, moves left, and $C$ reads $[y{-}2]$ --- halt iff $\tau(y{-}1)=0$ and $\tau(y{-}2)=0$: this is H2. Apart from $A0$ and $F1$ there is no other in-edge of $B$ reading a $1$ (row-by-row check of the table: the only transitions with target $B$ are $A0$ and $F1$). The converse directions construct the halting paths directly.
\end{proof}

\section{Sweep conservation and the lethal set}

\begin{definition}[empty-tail sweep]
A marker state with empty tail ($T=\varnothing$) is an empty-tail marker. Along eras, the $(L,M,R)$ of empty-tail $A$-markers evolves by Theorem~\ref{thm:dry}; such a trajectory is a \textbf{sweep}.
\end{definition}

\begin{theorem}[sweep conservation]\label{thm:sweep}
Along any sweep: $\mathrm{pop}=L{+}M{+}R$ is conserved; $Q:=L-3R$ is conserved; the terminal $M$ satisfies $M\equiv M_0\pmod 4$ ($M_0$ the chain-initial $M$). The reset maps $\Psi$ are
\[
\Psi_{M4}\colon(L,M,R)\mapsto(3,\,\mathrm{pop}-1,\,1),\qquad
\Psi_{M3}\colon\mapsto(L{+}8,\,R{-}4,\,1),\qquad
\Psi_{M2}\colon R_0'=R+2,\quad \Psi_{M1}\colon R_0'=1 .
\]
\end{theorem}
\begin{proof}
Conservation: $(L{+}3)+(M{-}4)+(R{+}1)=L{+}M{+}R$ and $(L{+}3)-3(R{+}1)=L-3R$; $M\equiv M_0\bmod4$ since each era applies $M\mapsto M-4$. The formulas for $\Psi$ are the parameter expansions of the corresponding terminal branches of Theorems~\ref{thm:D1}/\ref{thm:dispatch} (instance verification $164/164$, $162/166$, $171/171$ --- the $162/166$ of $\Psi_{M3}$ is the exact count on the $R\ge5$ domain, the 4 instances with $R\le4$ being dispatched separately by the symbolic branch D4, hence no mismatch).
\end{proof}

\begin{definition}[lethal set $P$]
$P$ is the set of marker parameter cells satisfying one of: (i) family g1g1 with $R=1$; (ii) family g3 with tail front $(g_1,1)$; (iii) family g5B with $L=2$; (iv) family g6o1 with $L=2$; (v) family g2L with $M=1$; (vi) family g6o2 with tail front $(g_1,1)$; (vii) family g4A with $R=1$ (via the g1g1-R2 prepend path). Its closure under preimages of the dispatch maps comprises \textbf{193 cells}.
\end{definition}

\begin{theorem}[lethality]\label{thm:lethal}
Every cell $p\in P$ is bare-machine-verified to halt: from the corresponding configuration the machine halts within 2--65 eras; in particular the synthetic instances $(L,M,R)=(7,2,1)$, $(8,2,2)$, $(21,1,1)$ halt after 33, 28, 24 eras respectively. Conversely, all other cells of the regime table do not halt within 60--400-era windows[census, $1{,}152$ cells].
\end{theorem}

\begin{theorem}[reduction]\label{thm:reduction}
The following are equivalent:
\begin{enumerate}
\item $M$, started on the blank tape, halts;
\item the era orbit $(\mathfrak M_k)_{k\ge0}$ of the initial marker under Theorems~\ref{thm:dry}/\ref{thm:D1}/\ref{thm:dispatch} satisfies $\exists k:\ \mathfrak M_k\in P$.
\end{enumerate}
Moreover, machine verification gives $\mathfrak M_k\notin P$ for all $k\le 2{,}830{,}885$ (i.e.\ $3\times10^{10}$ steps).
\end{theorem}
\begin{proof}
(2)$\Rightarrow$(1): Theorem~\ref{thm:lethal}. (1)$\Rightarrow$(2): suppose a halt; by Theorem~\ref{thm:H} the last state before halting realizes H1 or H2 inside some era. A per-family landing audit (families A/g5B/g6/g3/g2L/g4A/g4 have sprint-block right gap $\ge2$; family g1g1 with gap $=1$ has H1 $\iff R=1$; families g3/g6o2 have H1 $\iff$ tail front $(g_1,1)$; the blank-scan geometry of H2 occurs only for g5B-L2/g6o1-L2) shows H1/H2 trigger exactly when that era's marker state lies in $P$[landing audit + 68-cell halt census].
\end{proof}

\section{One entry sealed}

\begin{theorem}[M4 entry closed]\label{thm:m4}
If an empty-tail sweep is generated by an $M4$-reset ($\Psi_{M4}$), then its descendant M3-terminal $R$ values never enter the lethal lineage $\{8\cdot4^n-1:n\ge0\}=\{1,7,31,127,\dots\}$.
\end{theorem}
\begin{proof}
On a $\Psi_{M4}$-generated sweep, $L=3R$ holds exactly (Theorem~\ref{thm:sweep} with $Q=0$), hence $\mathrm{pop}=4R$; if the terminal $M=4$, the last era's output makes $\mathrm{pop}=4R{+}4$ --- in either case $\mathrm{pop}\equiv0\pmod4$, and the M3-terminal inherits $\mathrm{pop}\equiv0\pmod4$, while every lineage value satisfies $8\cdot4^n-1\equiv3\pmod4$. Contradiction.
\end{proof}

The initial marker, the g6 entries and the M3$\to$M3 self-closure of the lineage are checked branch-by-branch likewise[dispatch entry classification]. The unsealed entries are the M2- and M1-resets.

\section{The open problem}\label{sec:open}

\begin{conjecture}[trajectory avoidance]\label{conj:avoid}
The era orbit of the blank-tape initial marker satisfies $\forall k\ge0:\ \mathfrak M_k\notin P$.
\end{conjecture}

Verified: $k\le2{,}830{,}885$ ($3\times10^{10}$ steps); the minimum parameter distance from orbit-compatible cells to $P$ is $3$ (one instance at $t=34{,}922$: $(M1,L{=}48,R{=}20)$ versus $(M1,L{=}51,R{=}20)$), and the diagonal sweep motion ($\Delta L=3,\Delta R=1$ together) makes cells of different $Q$ unreachable in one step (Theorem~\ref{thm:sweep}). Refuted lines of attack:

\begin{enumerate}
\item \textbf{Route-constant exclusion is false}: the instances of Theorem~\ref{thm:lethal} show the M2/M1-reset route constants do not avoid the lineage preimage sets universally.
\item \textbf{No residue-class separation}: modulo $3\times3$, the reachable residue classes collide with all 193 cells of $P$; moreover the M2/M1-reset values $M_0'$ are not residue-class determined[an L6-type negative result].
\item \textbf{Linear invariants are insufficient}: conservation of $Q$ does not imply avoidance --- the configuration g6o1$(11,1,\varnothing)$ manufactures a tail front $(g_1,1)$ from an empty tail[engine reproduction]; on the real orbit $L\equiv0\pmod3$ renders $L\pm1$ cells unreachable (protecting 143/193 cells), but this invariant is an orbit observation without a written induction.
\end{enumerate}

\begin{remark}
Formal location: $(\mathfrak M_k)$ is the orbit of a piecewise-affine system; the avoidance problem belongs to the reachability class of piecewise-affine iterations (the family of concrete instances of Conway's 1972 undecidability \cite{conway72}; cf.\ the independence results of Kurtz--Simon 2007 \cite{kurtzsimon07} --- see paradigm-audit.md). ZFC-independence of this machine's halting cannot be excluded.
\end{remark}

\section*{Reproduction}
\texttt{ocm.exe} (bare-machine oracle); \texttt{xval-macro.mjs} (triple cross-validation); \texttt{census2/3.mjs} (regime census; the gE grid includes $M{=}3$); \texttt{tailfront.mjs} (tail-front evolution); \texttt{dispatch.mjs} (dispatch tables); \texttt{hclosure2.mjs} (family-graph fixed point, 294 nodes; its dependency \texttt{ref-lib.mjs} is vendored into this directory, one-line runnable); \texttt{shape2.mjs 1e10} (deep scan); \texttt{lethal-db.json} (the set $P$). All one-line commands; see \texttt{dossier.md}.

\begin{thebibliography}{9}
\bibitem{conway72} J.~H. Conway, \emph{Unpredictable iterations}, in: Proc. 1972 Number Theory Conference (Boulder), 49--52. University of Colorado, 1972.
\bibitem{kurtzsimon07} S.~A. Kurtz and J.~Simon, \emph{The undecidability of the generalized Collatz problem}, in: Theory and Applications of Models of Computation (TAMC 2007), LNCS 4484, 542--553. Springer, 2007.
\end{thebibliography}

\end{document}
